% =========================================================
% Vector Subspaces --- Notes
% =========================================================

\begin{remark*}[Toolkit --- Vector Subspaces]
\textbf{Concept family:} Subspaces, their sums, and direct sums.

\medskip
\textbf{Atomic items in this section:}
\begin{itemize}
  \item Definition: Subspace
  \item Theorem: Conditions for a Subspace
  \item Definition: Sum of Subspaces
  \item Theorem: Sum of Subspaces Is the Smallest Containing Subspace
  \item Definition: Direct Sum
  \item Theorem: Condition for a Direct Sum
  \item Theorem: Direct Sum of Two Subspaces
\end{itemize}

\medskip
\textbf{Key predicates:}
\[
  \operatorname{Subspace}(U,V),
  \qquad
  \operatorname{ClosedUnderAddition}(U),
  \qquad
  \operatorname{ClosedUnderScalarMultiplication}(U,\mathbf{F}),
  \qquad
  \operatorname{DirectSum}(V,\{U_1,\ldots,U_m\},\mathbf{F}).
\]
\textbf{Key test:} $U$ is a subspace $\iff$ $0\in U$, $U$ closed under $+$,
$U$ closed under scalar multiplication.
\end{remark*}

% =========================================================
% Definition --- Subspace
% =========================================================

\begin{definitionbox}{Definition (Subspace)}
\begin{definition}[Subspace]
\label{def:subspace}
Let $V$ be a vector space over $\mathbf{F}$.
A subset $U \subseteq V$ is called a \textbf{subspace} of $V$ if $U$ is itself
a vector space over $\mathbf{F}$ with the same addition and scalar multiplication
as on $V$.
\end{definition}
\end{definitionbox}

\begin{remark*}[Standard quantified statement]
\[
\begin{aligned}
&U \text{ is a subspace of } V \\
&\quad\Longleftrightarrow\quad
U \subseteq V
\text{ and } U \text{ is a vector space over } \mathbf{F} \\
&\qquad\text{under the operations inherited from } V.
\end{aligned}
\]
\end{remark*}

\begin{remark*}[Definition predicate reading]
\[
\operatorname{Subspace}(U,V)
\;\colonequiv\;
U \subseteq V \;\wedge\; \operatorname{VectorSpace}(U,\mathbf{F},+,\cdot).
\]
\end{remark*}

\begin{remark*}[Negated quantified statement]
\[
\neg\,\operatorname{Subspace}(U,V)
\;\Longleftrightarrow\;
U \not\subseteq V
\;\vee\;
\neg\,\operatorname{VectorSpace}(U,\mathbf{F},+,\cdot).
\]
\end{remark*}

\begin{remark*}[Interpretation]
A subspace is a subset of a vector space that is closed under the vector-space
operations and therefore forms a vector space in its own right.
\end{remark*}

\begin{dependencies}
\begin{itemize}
  \item \hyperref[def:vector-space]{Definition: Vector Space}
\end{itemize}
\end{dependencies}

% =========================================================
% Theorem --- Conditions for a Subspace
% =========================================================

\begin{theorembox}{Theorem (Conditions for a Subspace)}
\begin{theorem}[Conditions for a Subspace]
\label{thm:conditions-for-a-subspace}
Let $V$ be a vector space over $\mathbf{F}$ and let $U \subseteq V$.
Then $U$ is a subspace of $V$ if and only if the following three conditions hold:
\begin{enumerate}[label=(\roman*)]
  \item $0 \in U$;
  \item for all $u,v \in U$, one has $u+v \in U$;
  \item for all $a \in \mathbf{F}$ and all $u \in U$, one has $au \in U$.
\end{enumerate}
\hyperref[prf:conditions-for-a-subspace]{\textit{Go to proof.}}
\end{theorem}
\end{theorembox}

\begin{remark*}[Standard quantified statement]
\[
\begin{aligned}
&U \text{ is a subspace of } V \\
&\quad\Longleftrightarrow\quad
0 \in U \\
&\qquad\land\;
\operatorname{ClosedUnderAddition}(U) \\
&\qquad\land\;
\operatorname{ClosedUnderScalarMultiplication}(U,\mathbf{F}).
\end{aligned}
\]
\end{remark*}

\begin{remark*}[Failure modes]
$U$ fails to be a subspace via exactly one of three routes: it omits the zero
vector, it contains two vectors whose sum falls outside $U$, or it contains a
vector and a scalar whose product falls outside $U$.
\end{remark*}

\begin{remark*}[Contrapositive quantified statement]
\[
\neg\,\operatorname{Subspace}(U,V)
\;\Longleftrightarrow\;
0 \notin U
\;\vee\;
\neg\,\operatorname{ClosedUnderAddition}(U)
\;\vee\;
\neg\,\operatorname{ClosedUnderScalarMultiplication}(U,\mathbf{F}).
\]
\end{remark*}

\begin{remark*}[Interpretation]
To check that a subset is a subspace, it suffices to verify the three
conditions. The remaining seven vector space axioms are inherited automatically
from $V$ since $U$ uses the same operations. The three conditions are minimal:
condition (i) ensures $U$ is non-empty; (ii) and (iii) ensure the inherited
operations stay within $U$.
\end{remark*}

\begin{dependencies}
\begin{itemize}
  \item \hyperref[def:subspace]{Definition: Subspace}
  \item \hyperref[def:vector-space]{Definition: Vector Space}
\end{itemize}
\end{dependencies}

% =========================================================
% Definition --- Sum of Subspaces
% =========================================================

\begin{definitionbox}{Definition (Sum of Subspaces)}
\begin{definition}[Sum of Subspaces]
\label{def:sum-of-subspaces}
Let $V$ be a vector space over $\mathbf{F}$, and let
$U_1,\dots,U_m$ be subspaces of $V$.
The \textbf{sum} of $U_1,\dots,U_m$ is the set
\[
U_1+\cdots+U_m
=
\{u_1+\cdots+u_m : u_1 \in U_1,\dots,u_m \in U_m\}.
\]
\end{definition}
\end{definitionbox}

\begin{remark*}[Standard quantified statement]
\[
\begin{aligned}
&v \in U_1+\cdots+U_m \\
&\quad\Longleftrightarrow\quad
\exists u_1 \in U_1 \;\cdots\; \exists u_m \in U_m
\text{ such that } v=u_1+\cdots+u_m.
\end{aligned}
\]
\end{remark*}

\begin{remark*}[Definition predicate reading]
\[
\operatorname{SumOfSubspaces}(W,\{U_1,\dots,U_m\},V)
\;\colonequiv\;
W=U_1+\cdots+U_m.
\]
\end{remark*}

\begin{remark*}[Interpretation]
The sum of subspaces consists of all vectors that can be built by adding
together one vector from each subspace. Vectors from different subspaces
interact; the representation $v = u_1 + \cdots + u_m$ need not be unique.
\end{remark*}

\begin{dependencies}
\begin{itemize}
  \item \hyperref[def:subspace]{Definition: Subspace}
\end{itemize}
\end{dependencies}

% =========================================================
% Theorem --- Sum of Subspaces Is the Smallest Containing Subspace
% =========================================================

\begin{theorembox}{Theorem (Sum of Subspaces Is the Smallest Containing Subspace)}
\begin{theorem}[Sum of Subspaces Is the Smallest Containing Subspace]
\label{thm:sum-of-subspaces-smallest-containing}
Let $V$ be a vector space over $\mathbf{F}$, and let
$U_1,\dots,U_m$ be subspaces of $V$.
Then $U_1+\cdots+U_m$ is a subspace of $V$ containing each of
$U_1,\dots,U_m$.
Moreover, if $W$ is a subspace of $V$ containing each of
$U_1,\dots,U_m$, then
\[
U_1+\cdots+U_m \subseteq W.
\]
\hyperref[prf:sum-of-subspaces-smallest-containing]{\textit{Go to proof.}}
\end{theorem}
\end{theorembox}

\begin{remark*}[Standard quantified statement]
\[
\begin{aligned}
&\operatorname{Subspace}(U_1+\cdots+U_m,\,V), \\
&U_j \subseteq U_1+\cdots+U_m \quad \text{for } j=1,\dots,m, \\
&\forall W \subseteq V\;
\Bigl[
\bigl(W \text{ is a subspace of } V
\land U_1 \subseteq W \land \cdots \land U_m \subseteq W\bigr) \\
&\qquad\qquad\Rightarrow\;
U_1+\cdots+U_m \subseteq W
\Bigr].
\end{aligned}
\]
\end{remark*}

\begin{remark*}[Interpretation]
The sum is the smallest subspace containing all the given subspaces: it is
their join in the lattice of subspaces of $V$. Any subspace containing all
of $U_1,\ldots,U_m$ must contain every sum $u_1+\cdots+u_m$, so it
contains the sum subspace.
\end{remark*}

\begin{dependencies}
\begin{itemize}
  \item \hyperref[def:sum-of-subspaces]{Definition: Sum of Subspaces}
  \item \hyperref[thm:conditions-for-a-subspace]{Theorem: Conditions for a Subspace}
\end{itemize}
\end{dependencies}

% =========================================================
% Definition --- Direct Sum
% =========================================================

\begin{definitionbox}{Definition (Direct Sum)}
\begin{definition}[Direct Sum]
\label{def:direct-sum}
Let $V$ be a vector space over $\mathbf{F}$, and let
$U_1,\dots,U_m$ be subspaces of $V$ such that
\[
V=U_1+\cdots+U_m.
\]
Then $V$ is called the \textbf{direct sum} of $U_1,\dots,U_m$ if each
vector $v \in V$ can be written in the form
\[
v=u_1+\cdots+u_m,
\qquad
u_1 \in U_1,\dots,u_m \in U_m,
\]
in exactly one way.
In this case, we write
\[
V=U_1\oplus\cdots\oplus U_m.
\]
\end{definition}
\end{definitionbox}

\begin{remark*}[Standard quantified statement]
\[
\begin{aligned}
&V=U_1\oplus\cdots\oplus U_m \\
&\quad\Longleftrightarrow\quad
V=U_1+\cdots+U_m \\
&\qquad\land\;
\forall v \in V\;
\exists!\, u_1 \in U_1 \;\cdots\; \exists!\, u_m \in U_m
\text{ such that } v=u_1+\cdots+u_m.
\end{aligned}
\]
\end{remark*}

\begin{remark*}[Definition predicate reading]
\[
\operatorname{DirectSum}(V,\{U_1,\dots,U_m\},\mathbf{F})
\;\colonequiv\;
V=U_1\oplus\cdots\oplus U_m.
\]
\end{remark*}

\begin{remark*}[Negated quantified statement]
\[
\neg\,\operatorname{DirectSum}(V,\{U_1,\dots,U_m\},\mathbf{F})
\;\Longleftrightarrow\;
\exists v \in V\;\exists \text{ two distinct decompositions of } v.
\]
\end{remark*}

\begin{remark*}[Interpretation]
A direct sum is a sum in which every vector has a unique decomposition
into pieces coming from the given subspaces. The uniqueness condition is
the structural rigidity that makes direct sums so useful --- it means the
subspaces contribute independently, with no redundancy.
\end{remark*}

\begin{dependencies}
\begin{itemize}
  \item \hyperref[def:sum-of-subspaces]{Definition: Sum of Subspaces}
\end{itemize}
\end{dependencies}



\begin{tikzpicture}[scale=1.5]

% Draw the basis axes (R^3)
\draw[->, thin, gray] (-1,0,0) -- (4,0,0) node[right] {$x$};
\draw[->, thin, gray] (0,-1,0) -- (0,3,0) node[above] {$y$};
\draw[->, thin, gray] (0,0,-1) -- (0,0,3) node[below left] {$z$};

% Define the Subspaces U1 and U2
% U1 is a plane in R^3
\coordinate (U1_top) at (3.5,2.5,0.5);
\coordinate (U1_bot) at (-0.5,-0.5,1.5);
\fill[blue!10, opacity=0.7] (0,0,0) -- (3,2,0) -- (3.5,2.5,1) -- (0.5,0.5,1) -- cycle;
\draw[blue!50, dashed] (0,0,0) -- (3,2,0);
\draw[blue!50, dashed] (0,0,0) -- (0.5,0.5,1);
\node[blue] at (1.5,1.2,0.8) {$U_1$};

% U2 is a line in R^3
\coordinate (U2_origin) at (0,0,0);
\coordinate (U2_end) at (0.5,-1,3);
\draw[red, thick] (U2_origin) -- (U2_end);
\draw[red, thin, dashed] (0,0,0) -- (-0.2,0.4,-1.2); % Extension below origin
\node[red] at (0.3,-0.7,2) {$U_2$};

% Define vector v and its unique decomposition
\coordinate (O) at (0,0,0);
\coordinate (u1) at (2, 1.33, 0.67); % A vector strictly in U1 (blue plane)
\coordinate (u2) at (0.2, -0.4, 1.2); % A vector strictly in U2 (red line)
\coordinate (v) at ($(u1)+(u2)$); % Vector v = u1 + u2

% Draw Vector v (The sum)
\draw[->, ultra thick, black, line cap=round] (O) -- (v) node[above right] {$\mathbf{v} \in V$};

% Draw the unique components (The basis of the direct sum)
\draw[->, thick, blue!80, line cap=round] (O) -- (u1) node[midway, above left, xshift=-2pt] {$\mathbf{u}_1 \in U_1$};
\draw[->, thick, red!80, line cap=round] (O) -- (u2) node[midway, left] {$\mathbf{u}_2 \in U_2$};

% Draw projection lines showing the unique parallelogram (Decomposition)
\draw[thin, dashed, gray] (u1) -- (v);
\draw[thin, dashed, gray] (u2) -- (v);

% Highlight the unique point of intersection (The zero vector)
\fill[black] (O) circle (2pt);
\node[below, xshift=3pt] at (O) {$\mathbf{0}$};

% Label the overall space and uniqueness condition
\node[anchor=north, align=center, text width=5cm] at (1.5, -1.5) {
    $V = U_1 \oplus U_2$ \\
    \textbf{Uniqueness}: $\mathbf{v} = \mathbf{u}_1 + \mathbf{u}_2$ \\
    \textbf{Trivial Intersection}: $U_1 \cap U_2 = \{\mathbf{0}\}$
};

\end{tikzpicture}



% =========================================================
% Theorem --- Condition for a Direct Sum
% =========================================================

\begin{theorembox}{Theorem (Condition for a Direct Sum)}
\begin{theorem}[Condition for a Direct Sum]
\label{thm:condition-for-direct-sum}
Let $V$ be a vector space over $\mathbf{F}$, and let
$U_1,\dots,U_m$ be subspaces of $V$ such that
\[
V=U_1+\cdots+U_m.
\]
Then the following are equivalent:
\begin{enumerate}[label=(\roman*)]
  \item \(V=U_1\oplus\cdots\oplus U_m\);
  \item if \(u_1 \in U_1,\dots,u_m \in U_m\) and
        \[
        u_1+\cdots+u_m=0,
        \]
        then
        \[
        u_1=\cdots=u_m=0.
        \]
\end{enumerate}
\hyperref[prf:condition-for-direct-sum]{\textit{Go to proof.}}
\end{theorem}
\end{theorembox}

\begin{remark*}[Standard quantified statement]
\[
\begin{aligned}
&V=U_1\oplus\cdots\oplus U_m \\
&\quad\Longleftrightarrow\quad
\forall u_1 \in U_1 \;\cdots\; \forall u_m \in U_m \\
&\qquad
\bigl(u_1+\cdots+u_m=0 \Rightarrow u_1=\cdots=u_m=0\bigr).
\end{aligned}
\]
\end{remark*}

\begin{remark*}[Interpretation]
A sum is direct exactly when the zero vector has only the trivial
decomposition as a sum of vectors from the given subspaces. This is the
practical test: rather than verifying uniqueness for every $v \in V$,
it suffices to check uniqueness at $v = 0$.
\end{remark*}

\begin{dependencies}
\begin{itemize}
  \item \hyperref[def:direct-sum]{Definition: Direct Sum}
  \item \hyperref[thm:unique-additive-identity]{Theorem: Unique Additive Identity}
\end{itemize}
\end{dependencies}

% =========================================================
% Theorem --- Direct Sum of Two Subspaces
% =========================================================

\begin{theorembox}{Theorem (Direct Sum of Two Subspaces)}
\begin{theorem}[Direct Sum of Two Subspaces]
\label{thm:direct-sum-of-two-subspaces}
Let $U$ and $W$ be subspaces of a vector space $V$.
Then
\[
U+W=U\oplus W
\]
if and only if
\[
U \cap W = \{0\}.
\]
\hyperref[prf:direct-sum-of-two-subspaces]{\textit{Go to proof.}}
\end{theorem}
\end{theorembox}

\begin{remark*}[Standard quantified statement]
\[
\begin{aligned}
&U+W=U\oplus W \\
&\quad\Longleftrightarrow\quad
U \cap W = \{0\}.
\end{aligned}
\]
\end{remark*}

\begin{remark*}[Contrapositive quantified statement]
\[
U \cap W \neq \{0\}
\;\Longleftrightarrow\;
U + W \neq U \oplus W.
\]
If the two subspaces share a nonzero vector, the sum is not direct.
\end{remark*}

\begin{remark*}[Interpretation]
For two subspaces, directness is equivalent to saying they overlap only
in the zero vector. This gives a simple geometric test: draw both subspaces
and check whether their intersection is just the origin. The result does not
generalise to three or more subspaces --- it is possible for $U_1\cap U_2$,
$U_1\cap U_3$, and $U_2\cap U_3$ to all equal $\{0\}$ yet $U_1+U_2+U_3$
still fail to be a direct sum.
\end{remark*}

\begin{dependencies}
\begin{itemize}
  \item \hyperref[thm:condition-for-direct-sum]{Theorem: Condition for a Direct Sum}
  \item \hyperref[def:direct-sum]{Definition: Direct Sum}
\end{itemize}
\end{dependencies}
